This thesis focuses on the use of reinforcement learning for adversarial planning within the framework of realistic guidance, navigation, and controls for orbital spacecraft. Traditional approaches based on classical optimization methods can become brittle when faced against intelligent and non-cooperative opponents which can have shifting strategies. Reinforcement learning algorithms can be useful tools to fill in this gap, so long as they're used within the right use case and context that mitigates these algorithms' lack of convergence guarantees, proneness to instability, non-stationarity, and limited utilization outside the training dataset.

With this in mind, this thesis proposes the Phased Adversarial Reinforcement Learning plus Behavior Cloning (\acs{parlbc}, for short) framework to train robust policies in a two-player, zero-sum Target-Attacker-Defender (\acs{tad}) orbital defense scenario. This framework is built based on the well-known Proximal Policy Optimization (\acs{ppo}) online reinforcement learning (\acs{rl}) algorithm. We train attackers by freezing existing opponent policies and using randomized mixtures of these, with the intent of mitigating non-stationarity during learning and avoiding the instabilities that are often seen on simultaneous play methods. 

Overall, the overarching goal of the thesis and the \acs{parlbc} framework as a whole to develop a phased \acs{ppo}-based learning framework that makes it possible to apply \acs{rl} formulations in a dynamic adversarial problem with physical constraints. To test this, we study a two-player zero-sum orbital adversarial planning, extending the framework to partial observability of an opponent's state through behavior cloning with bearings-only Kalman-filter state estimates, studying how learned policies interact with a maximum velocity limiting classical \acs{cbfqp} safety controller in-the-loop, and evaluating if policies trained under \acs{hcw} dynamics can transfer to previously unseen elliptical linear time-varying dynamics.


These learned policies train within a closed-loop aerospace guidance, navigation, and control architecture (\acs{gnc}). The control piece comes in the form of a Control Barrier Function-Quadratic Program (\acs{cbfqp}), used to enforce the feasibility of a learned control action subject to a vehicle's maximum velocity constrained relative to a predefined safety boundary. Moreover, the navigation piece, which is meant to showcase the feasibility of the framework under partial observability simulating real-life sensor limitations, is addressed by doing behavior cloning from observations obtained from an Extended Kalman Filter (\acs{ekf}) in the loop.

We trained three different test cases that resulted in multiple attacker and defender policies across consecutive matchups, which showed stability during learning and convergent behavior.

Furthermore, the policies are evaluated via Monte Carlo rollouts under both full and partial observability, varying spherical arena sizes (e.g: with different radii), and the proper dispersions to randomize the starting states of each agent as well as their \acs{ekf} parameters. Although policies are trained under Hill-Clohessy-Wiltshire (\acs{hcw}) linear time-invariant relative orbital dynamics, they are also evaluated under elliptical dynamics to test if the learned policies can transfer to dynamics not seen during training.

The Monte Carlo results show that our proposed framework produces coherent and competitive defender and attacker policies for a zero-sum orbital defense game in which a defender spacecraft must prevent the attacker from reaching a protected asset by intercepting it first within a bounded spherical arena. The agents are assumed to undergo relative motion within a spherical arena centered on the chief of the orbit, which is the asset being defended, while the entire arena orbits the Earth with the chief. Within a small, 20 m training spherical arena, full-state policies reached up to a 79.9\% defender win rate against a competitive \acs{rl} attacker with access to the opponent's full state, and 74.7\% for the \acs{ekf}-based policies with bearings-only measurements of the opponent's state, a 6.51\% drop in performance relative with respect to\ the fully observable policies.

When evaluating in a larger 100m radius spherical arena, not used or seen by the policies during training, full-state results achieved up to a 66.2\% defender win rate, a 17.15\% relative drop in performance from the 20m radius training arena, whereas \acs{ekf}-based policies struggled more with a 47.4\% defender win rate, a 28.4\% relative drop from the full state case. 

Evaluating all of the above across both \acs{hcw} and elliptic dynamics, even though the policies were never exposed to the latter during training, resulted in at most a 1.29\% relative change between results. Therefore, within the cases tested in this work, the policies trained under dynamics transferred effectively to unmodeled dynamics used during evaluation.

Throughout the Monte Carlo runs, median policy inference time reached up to 0.145 ms per timestep on the full-state policies and 3.91 ms per timestep on the \acs{ekf}-based policies, the latter of which includes the \acs{ekf} compute time in the loop.

All in all, the results obtained via our proposed PARL+BC method and its evaluation show that reinforcement learning can be integrated with classical navigation and control architectures in the loop in constrained adversarial orbital scenarios, resulting in robust and computationally efficient policies that work under partial observability and can retain their performance over mild orbital perturbations and dynamics conditions beyond those seen during training.
